% TOP_PROBLEM: {"id": 6200044, "title": "Boundaries of Groups and Kleinian Groups — Problem 44", "category_id": 6, "category": {"id": 6, "name": "geometry", "display_name": "Geometry", "description": "Euclidean and non-Euclidean geometry, geometric structures.", "slug": "geometry", "order_index": 6, "created_at": "2026-07-31T15:26:25.671Z"}, "status": "open", "problem_number": "AMR-061-0044", "set_id": 13, "statusQualification": "Natural nontrivial measure class is made precise in the analytic setting as a nonatomic full-support Radon class; invariance means mutual absolute continuity."}
% FIRST_POSTED: 2026-10-02
% ENTRY_KIND: statement_only
% UnsolvedMath ID 6200044; source catalogue code AMR-061-0044
% !TeX program = lualatex
% Definition and sources only; this file is not an LLM solution attempt.
\documentclass[11pt,a4paper]{article}
\usepackage[margin=25mm]{geometry}
\usepackage{fontspec}
\setmainfont{lmroman10-regular.otf}[BoldFont=lmroman10-bold.otf,
 ItalicFont=lmroman10-italic.otf,BoldItalicFont=lmroman10-bolditalic.otf]
\usepackage{amsmath,amssymb,mathtools}
\usepackage{unicode-math}
\setmathfont{latinmodern-math.otf}
\AtBeginDocument{\let\setminus\smallsetminus}
\usepackage{xurl,hyperref}
\hypersetup{colorlinks=true,urlcolor=blue}
\setcounter{secnumdepth}{0}
\setlength{\parindent}{0pt}
\setlength{\parskip}{5pt}
\setlength{\emergencystretch}{3em}
\begin{document}
\section*{Boundaries of Groups and Kleinian Groups — Problem 44}\label{6200044}
{\small UnsolvedMath ID 6200044 · AMR-061-0044 · Geometry · AMR Open Problem Lists}
\subsection{Definitions and mathematical statement}
Let $G$ be a word-hyperbolic group whose Gromov boundary $Z=\partial G$ is connected and has no local cut points. A local cut point is a point whose removal disconnects some connected open neighborhood. Equip $Z$ with a visual metric, defined up to quasisymmetric equivalence by the lengths of common initial segments of geodesic rays.

A homeomorphism $f:Z\to Z$ is quasisymmetric if a homeomorphism $\eta:[0,\infty)\to[0,\infty)$ controls every distance ratio by
\[
 \frac{d(fx,fa)}{d(fx,fb)}\leq
                 \eta\!\left(\frac{d(x,a)}{d(x,b)}\right).
\]
A measure class $[\mu]$ consists of the Borel measures with exactly the same null sets as $\mu$. It is invariant under quasisymmetric homeomorphisms when, for every such $f$,
\[
 f_*\mu\sim\mu,
\]
where $\sim$ denotes mutual absolute continuity.

Can one construct a natural, geometrically defined nontrivial measure class on $Z$ with this invariance? In the intended analytic formulation seek a nonatomic full-support Radon measure, normalized to probability on the compact boundary. Thus single points have measure zero and every nonempty open set has positive measure.

Invariance concerns null sets, not exact preservation of the numerical measure of every set. It is required for the full quasisymmetric homeomorphism group, not only the boundary action of $G$. Merely choosing a Patterson--Sullivan measure without proving the additional absolute-continuity statement does not answer the question outside known special cases.
\subsection{Short English statement}
Does every connected hyperbolic boundary without local cut points have a natural measure class preserved by all quasisymmetries?
\subsection{Sources}
\begin{itemize}
\item[S1] ULAM.AI and the UnsolvedMath Contributors, supplied problems.json, record 6200044 (AMR-061-0044), AMR Open Problem Lists. Catalogue identity and source transcription; CC BY 4.0.
\url{https://huggingface.co/datasets/ulamai/UnsolvedMath}.
\item[S2] Kapovich - Problems on Boundaries of Groups and Kleinian Groups (2005), Problem 44, PDF page 13. Original source indexed by AMR-061-0044.
\url{https://www.math.ucdavis.edu/~kapovich/EPR/problems.pdf}.
\end{itemize}
\end{document}
